\section*{Chapitre \ref{ch:b3:partition-functions} --- Thermodynamique statistique~: fonctions de partition}

\begin{solution}{exo:b3:partition-functions:1}
Notons chaque distribution par les nombres de quanta des quatre
oscillateurs, rangés par ordre décroissant. «~4, 0, 0, 0~»~: $4!/(3!\,1!) = 4$~; «~3, 1, 0, 0~»~: $4!/(2!\,1!\,1!) = 12$~; «~2, 2, 0, 0~»~:
$4!/(2!\,2!) = 6$~; «~2, 1, 1, 0~»~: $4!/(1!\,2!\,1!) = 12$~; «~1, 1, 1, 1~»~: 1. Total 35 $= \binom73$.
Les plus probables sont «~3, 1, 0, 0~» et «~2, 1, 1, 0~» (12 chacune)~;
la seconde, avec des \omterm{def:b3:partition-functions:boltzmann}{populations} 1, 2, 1 pour 0, 1, 2 quanta, décroît
déjà avec l'énergie comme le fait la \omterm{def:b3:partition-functions:boltzmann}{distribution de Boltzmann} pour de
grands nombres.
\end{solution}

\begin{solution}{exo:b3:partition-functions:2}
$\eu^{-2500/(8.314 \times 298.15)} = \eu^{-1.009} = 0.365$~; à \qty{1000}{K}, $\eu^{-0.301} = 0.740$.
À très haute température, le rapport tend vers 1 (\omterm{def:b3:partition-functions:boltzmann}{populations} égales,
jamais d'inversion).
\end{solution}

\begin{solution}{exo:b3:partition-functions:3}
$m = 39.95\times10^{-3}/6.022\times10^{23} = \qty{6.634e-26}{kg}$~; $\Lambda = h/\sqrt{2\pi mkT} =
\qty{1.600e-11}{m} = \qty{16.0}{pm}$; $q^{\mathrm T} = 10^{-3}/(1.600\times10^{-11})^3 = \num{2.44e29}$.
\end{solution}

\begin{solution}{exo:b3:partition-functions:4}
$hc/k = \qty{1.4388}{cm.K}$. \ce{N2}~: $\theta_{\mathrm R} = \qty{2.863}{K}$, $q^{\mathrm R} = 298.15/(2 \times 2.863) = 52.1$.
\ce{HCl}~: $\theta_{\mathrm R} = \qty{15.02}{K}$, $q^{\mathrm R} = 298.15/15.02 = 19.85$
(somme exacte 20.19).
\end{solution}

\begin{solution}{exo:b3:partition-functions:5}
\ce{I2}~: $\theta_{\mathrm V} = 1.4388 \times 213.27 = \qty{306.9}{K}$, $q^{\mathrm V} = 1/(1 - \eu^{-1.029}) = 1.556$~;
fraction excitée $1 - 1/q^{\mathrm V} = 0.357$. \ce{N2}~: $\theta_{\mathrm V} = \qty{3352}{K}$, $q^{\mathrm V} = 1.000013$,
fraction excitée \num{1.3e-5}.
\end{solution}

\begin{solution}{exo:b3:partition-functions:6}
\ce{NO}~: $q^{\mathrm E} = 2 + 2\eu^{-119.73/207.22} = 2 + 2 \times 0.561 = 3.12$~;
fraction dans le niveau supérieur $1.12/3.12
= 0.36$. À haute température, $q^{\mathrm E} \to 4$ (deux niveaux
également peuplés). \ce{Cl}~:
$q^{\mathrm E} = 4 + 2\eu^{-882.35/207.22} = 4 + 2 \times 0.0142 = 4.03$.
\end{solution}

\begin{solution}{exo:b3:partition-functions:7}
$\ln q = \ln V + \frac32\ln T + \text{cste}$, donc $U - U(0) = NkT^2 \times \frac{3}{2T} = \frac32NkT$~;
pour une mole, $\frac32RT = \qty{3.72}{kJ/mol}$ à \qty{298.15}{K}, et $C_V = \frac32R = \qty{12.47}{J.K^{-1}.mol^{-1}}$.
\end{solution}

\begin{solution}{exo:b3:partition-functions:8}
$kT/p^\circ = \qty{4.116e-26}{m^3}$, $\Lambda^3 = \qty{4.093e-33}{m^3}$, rapport \num{1.006e7}~:
$S_{\mathrm m} = 8.3145 \times (16.124 + 2.5) = \qty{154.85}{J.K^{-1}.mol^{-1}}$,
la valeur tabulée au dernier chiffre près. Diviser $p$ par 10 ajoute
$R\ln10 = \qty{19.14}{J.K^{-1}.mol^{-1}}$.
\end{solution}

\begin{solution}{exo:b3:partition-functions:9}
$S^{\mathrm R}_{\mathrm m} = R[\ln(298.15/15.02) + 1] = 8.3145 \times (2.988 + 1) = \qty{33.16}{J.K^{-1}.mol^{-1}}$.
\end{solution}

\begin{solution}{exo:b3:partition-functions:10}
$\frac{\dd}{\dd J}\bigl[(2J+1)\eu^{-\theta J(J+1)/T}\bigr] = \bigl[2 - (2J+1)^2\theta/T\bigr]\eu^{-\theta J(J+1)/T} = 0$ donne
$2J + 1 = \sqrt{2T/\theta}$, soit $J_{\max} = \sqrt{T/2\theta} - \frac12$. \ce{HCl}~: 2.66, donc $J = 3$~; \ce{CO}~:
6.86, donc $J = 7$. L'intensité d'une raie de la \omterm{def:b3:rovibrational-spectroscopy:branches}{branche P} ou R suit la
\omterm{def:b3:partition-functions:boltzmann}{population} de son niveau inférieur, maximale en $J_{\max}$, à quelques
raies de la lacune du centre de la bande.
\end{solution}

\begin{solution}{exo:b3:partition-functions:11}
$\int_0^\infty f\dd J = T/\theta$, $f(0) = 1$, $f'(0) = 2 - \theta/T$. Donc $q^{\mathrm R} \approx T/\theta + \frac12 -
\frac16 + \frac{\theta}{12T} \approx T/\theta + \frac13$. L'erreur
relative de $T/\theta$ vaut environ $\theta/3T$~:
moins de 1\,\% quand $T > 33\,\theta$. Pour \ce{H2}, $\theta = 1.4388 \times 59.32 = \qty{85.3}{K}$~:
à \qty{298.15}{K}, l'erreur est de 10\,\%, inacceptable (et la
restriction sur $J$ due au spin nucléaire doit aussi être traitée
exactement, \cref{ch:b3:statistical-thermo-applied}).
\end{solution}

\begin{solution}{exo:b3:partition-functions:12}
Pour $x = \theta_{\mathrm V}/T \to \infty$, $x/(\eu^x - 1) \to 0$ et $\ln(1 - \eu^{-x}) \to 0$~: $S \to 0$. Pour $x \to 0$,
$x/(\eu^x - 1) \approx 1 - x/2$ et $-\ln(1 - \eu^{-x}) \approx -\ln x + x/2$~: $S \to R(1 - \ln x) = R[1 + \ln(T/\theta_{\mathrm V})]$.
\ce{I2}, $x = 1.029$~: valeur exacte $R(0.5721 + 0.4420) = \qty{8.43}{J.K^{-1}.mol^{-1}}$~;
limite $R(1 - 0.0289) =
\qty{8.07}{J.K^{-1}.mol^{-1}}$. Même pour $T \approx \theta_{\mathrm V}$,
la limite classique n'est trop basse que de 4\,\%.
\end{solution}

\begin{solution}{pb:b3:partition-functions:1}
\textbf{1.} $m = 28.014\times10^{-3}/6.02214\times10^{23} = \qty{4.652e-26}{kg}$.
\textbf{2.} $\Lambda = h/\sqrt{2\pi mkT} = \qty{1.910e-11}{m}$ (\qty{19.1}{pm}).
\textbf{3.} $kT/p^\circ = 1.380649\times10^{-23} \times 298.15/10^5 = \qty{4.116e-26}{m^3}$.
\textbf{4.} $q^{\mathrm T}/N = 4.116\times10^{-26}/(1.910\times10^{-11})^3 = \num{5.905e6}$~:
environ six millions d'états de translation par molécule, si bien que
deux molécules n'en partagent presque jamais un et que
$Q = q^N/N!$ est valable.
\textbf{5.} $S^{\mathrm T}_{\mathrm m} = R(\ln\num{5.905e6} + \frac52) = 8.3145 \times 18.091 = \qty{150.42}{J.K^{-1}.mol^{-1}}$.
\textbf{6.} $+R\ln2 = \qty{5.76}{J.K^{-1}.mol^{-1}}$, soit \num{156.18}.
\textbf{7.} $B_0 = 1.998241 - 0.008659 = \qty{1.989582}{cm^{-1}}$.
\textbf{8.} $\theta_{\mathrm R} = 1.438777 \times 1.989582 = \qty{2.8626}{K}$.
\textbf{9.} $\sigma = 2$~: les deux noyaux sont identiques~; retourner la
molécule bout pour bout donne une orientation indiscernable, et la
symétrie d'échange des noyaux n'autorise, pour chaque état de spin
nucléaire, que la moitié des niveaux de rotation.
\textbf{10.} $q^{\mathrm R} = 298.15/(2 \times 2.8626) = 52.08$.
\textbf{11.} Oui~: $T/\theta_{\mathrm R} = 104$, soit une erreur d'environ
0.3\,\% sur $q^{\mathrm R}$ et bien moindre sur l'entropie.
\textbf{12.} $S^{\mathrm R}_{\mathrm m} = R(\ln52.08 + 1) = 8.3145 \times 4.9527 = \qty{41.18}{J.K^{-1}.mol^{-1}}$.
\textbf{13.} $J_{\max} = \sqrt{298.15/5.725} - 0.5 = 6.72$~: $J = 7$
(\omterm{def:b3:partition-functions:boltzmann}{population} 0.0839, juste au-dessus de
0.0831 pour $J = 6$).
\textbf{14.} $2358.57 - 2 \times 14.324 = \qty{2329.92}{cm^{-1}}$.
\textbf{15.} $\theta_{\mathrm V} = 1.438777 \times 2329.92 = \qty{3352}{K}$.
\textbf{16.} $x = 3352.2/298.15 = 11.24$~; $q^{\mathrm V} = 1/(1 - \eu^{-11.24}) = 1.000013$.
\textbf{17.} $\eu^{-11.24}/q^{\mathrm V} = \num{1.3e-5}$.
\textbf{18.} $S^{\mathrm V}_{\mathrm m} = R[11.24\eu^{-11.24} + \eu^{-11.24}] = \qty{0.0013}{J.K^{-1}.mol^{-1}}$,
négligeable.
\textbf{19.} $g_0 = 1$ (${}^1\Sigma$)~: $R\ln1 = 0$. Les états
électroniques excités se situent plusieurs électronvolts plus haut, à
plus de cent fois $kT$~: leurs facteurs de Boltzmann sont totalement
négligeables.
\textbf{20.} $150.42 + 41.18 + 0.00 + 0 = \qty{191.60}{J.K^{-1}.mol^{-1}}$.
\textbf{21.} L'écart vaut \qty{0.01}{J.K^{-1}.mol^{-1}}, au niveau des
approximations faites (\omterm{def:b3:quantum-model-systems:rotor}{rotateur rigide}, séparation des mouvements).
\textbf{22.} Les capacités thermiques du solide mesurées depuis quelques
kelvins (extrapolées en $T^3$ au-dessous), les enthalpies de la
transition solide--solide, de fusion et de vaporisation divisées par
leurs températures, la capacité thermique du gaz, et une correction de
non-idéalité.
\textbf{23.} Les noyaux sont conservés dans toute réaction, si bien que
leur entropie de spin est la même des deux côtés et se simplifie~; la
calorimétrie ne la voit pas non plus, puisque les spins nucléaires
restent désordonnés dans le cristal jusqu'aux plus basses températures
mesurées.
\textbf{24.} \textbf{$S^\circ(\ce{N2}, \qty{298.15}{K}) = \qty{191.6}{J.K^{-1}.mol^{-1}}$
à partir des constantes spectroscopiques}, dont 150.4 dus à la
translation et 41.2 à la rotation.
\end{solution}
